% !TeX root = ../../main.tex
% cbic/errata.tex -- one table. Extraction provenance and verification: tables/README.md
% [round6] CONVERTED table -> longtable, for the same reason and by the same precedent as
% tables/frame/bib_errata.tex: with the Standley claim-narrowing row added, this table is TALLER THAN
% THE TEXT BLOCK, and a float cannot break across pages. Measured before converting: pdflatex reported
% "Float too large for page by 190.30908pt", the only oversized float in either build.
% Measured after converting: oversized_floats=0 in both variants, at 109 pages defense and 105 final.
% The conversion itself costs ZERO pages, isolated rather than assumed: in a scratch copy, this round's
% eight rows put back into the original float container build to the same 109/105, with the "Float too
% large" warning present in both variants. So the conversion removes the warning and moves no page.
% ATTRIBUTION OF THE PAGE GROWTH, measured and NOT what an earlier version of this note claimed. In a
% second scratch copy, with only this track's six files reverted to HEAD and every other file left as it
% currently stands, the build is 108/104 against the current 109/105. So this track's prose additions
% across Chapters 3, 4 and 5 and this appendix cost about ONE page per variant. The rest of the distance
% from the 106/101 that a pristine HEAD checkout produces belongs to other work landing in this tree in
% parallel, not to this track. Do not read 106 -> 109 as this round's cost; the isolated figure is
% 108 -> 109 defense and 104 -> 105 final.
% All eight rows, the caption, the label and \small are preserved; only the container changed.
% The {\small ... } GROUP is load-bearing and must stay: longtable is not a float, so \small cannot be
% set inside the environment the way a \begin{table} sets it. Dropping its opening brace is what stopped
% every build for six commits in round 5 (see the note in frame/bib_errata.tex). Brace balance verified
% after the change.
% \centering is dropped, not lost: longtable centers itself through \LTleft/\LTright, and a \centering
% inside it has no effect. Column widths, header text and row order are untouched.
{\small
\begin{longtable}{@{}p{0.42\textwidth}p{0.52\textwidth}@{}}
\caption{Content errata of the CBIC 2025 article and the corrections applied in
Chapter~\ref{ch:cbic}. Numeric values quoted in the corrections are the published
table values.}
\label{tab:apx:cbic-errata} \\
\toprule
\textbf{Defect in the published article} & \textbf{Correction in the chapter} \\
\midrule
\endfirsthead
\multicolumn{2}{@{}l}{\footnotesize\itshape Table \thetable\ continued from the previous page.}\\
\toprule
\textbf{Defect in the published article} & \textbf{Correction in the chapter} \\
\midrule
\endhead
\midrule
\multicolumn{2}{r@{}}{\footnotesize\itshape Continued on the next page.}\\
\endfoot
\bottomrule
\endlastfoot
Prose states ``almost four times more wall time''; the convergence table gives
80.88~s for the joint model against a cumulative 34.97~s for the single-task models,
a factor of about 2.3. &
Prose reconciled to the table: ``about 2.3 times the cumulative 34.97~s of the
individual single-task models''. \\
\addlinespace
Prose states the MFLOPs cost of the joint model is ``roughly double''; the
table gives 0.234 MFLOPs for the joint model against 2.315 and 0.012 for the two
single-task models. &
Prose reconciled to the table: the text now states that the table does not show a
higher MFLOPs cost for the joint model and quotes the three cell values; four
follow-on wording edits restrict the overhead claim to wall time. \\
\addlinespace
Broken cross-reference: the label intended for the single-task heads sits on the
Dataset subsection, and the method section points to it as if it described the heads. &
Dataset subsection relabeled; the dangling pointer sentence was removed, since no
section describing the single-task heads exists in the published paper. \\
\addlinespace
Typo ``spatio-tegm mporal''. &
``spatio-temporal''. \\
\addlinespace
Unfilled dataset placeholders ($N_{\text{users}}$, $N_{\text{poi}}$,
$N_{\text{checkins}}$) in the results section. &
Filled with the Florida figures of record for the same Gowalla data, as published in the
CoUrb dataset table: 20{,}301 users, 65{,}009 POIs, and 990{,}518 check-ins. The published
CBIC paper left these three statistics as placeholders; they are taken from the published
CoUrb table so that one corpus figure serves both chapters, and they are the counts that the
extraction pipeline of that time produced, not the counts the current pipeline produces
(Section~\ref{apx:errata:corpus}). \\
\addlinespace
The optimizer subsection states that Nash-MTL's ``reliance on gradient signs rather than
scales inherently balances heterogeneous tasks''. Nash-MTL does not read gradient signs,
and the same subsection states the correct property four paragraphs earlier, listing
scale invariance among the optimizer's axioms. &
Corrected to name the property the method has: its invariance to the scale of the
individual task gradients, the axiom already listed in the same subsection, balances
heterogeneous tasks without manual re-weighting. Checked against the optimizer's own
paper and against the implementation released with the two studies, neither of which
uses gradient signs. The claim is weaker than the published one, and no result depends
on it. \\
\addlinespace
Prose states that Nash-MTL ``requires only two matrix-vector products per iteration''. &
The clause is removed from the reproduced sentence and its correction given in a footnote.
The claim is not supported by the method's own paper, where the phrase does not occur, and
it understates the cost: both implementations examined run an iterative concave-convex
procedure, twenty passes by default, each a convex solve, on top of one backward pass per
task. The correction raises the reported cost of a method this chapter used, so it runs
against the chapter's own interest. \\
\addlinespace
The rationale for hard parameter sharing states that the approach ``frequently matches or
exceeds the performance of more complex architectures on many benchmarks, while offering
faster training and inference'', citing a work on task grouping. &
Narrowed to the inference-cost claim the cited work supports, that a single network is
evaluated rather than one network per task, with the correction recorded in a footnote. The
cited work names improved accuracy and reduced training time among the benefits joint
training may have in theory, then reports empirically that joint training often performs
worse than separate networks, and its contribution is a framework for splitting tasks across
networks, which it describes as costly at training time. Its one result favoring joint
training is conditioned on a matched parameter budget and on three or more tasks, and it
excludes the two-task case this chapter studies. The correction removes a stated advantage of
the architecture this chapter adopts, so it runs against the chapter's own interest. \\
\addlinespace
The survey of optimization techniques states that MGDA ``finds Pareto-optimal descent
directions''. Pareto optimality and a per-step descent direction are two different
guarantees, and the cited work proves the first for a bound rather than for the step. &
Narrowed to the two properties the cited work states. Its method computes a direction that
decreases every task loss, or else the current point is already Pareto-stationary; the Pareto
optimality it proves is for an upper bound on the multi-objective loss and holds under stated
assumptions. Neither the chapter's design nor any of its results depends on this method, which
it does not use. \\
\addlinespace
The rationale for hard parameter sharing attributes the claim that constraining the
hypothesis space regularizes, and generalizes better when tasks are related, to a work on
learned soft sharing. That work is a method for learning what and how much to share,
presented as an improvement on fixed sharing, so it argues against the sentence rather than
for it. &
Re-attributed to the work that states the argument the sentence makes: choosing a hypothesis
space large enough to contain a solution yet small enough to generalize, searched over an
environment of related tasks. The substituted reference was already in the bibliography and
is cited for the same claim elsewhere in the chapter. The soft-sharing work remains cited
where the chapter discusses cross-connection architectures and negative transfer. \\
\end{longtable}
}
